% =====================================================================
%  Banco complementar --- Fontes Reais e Leis de Kirchhoff  (REVISADO)
%  Substitui a versao gerada por template (9 clones em N2, 9 em N3 e
%  7 questoes conceituais marcadas como N4).
%  O cap02 ja cobre: duas fontes ideais em oposicao, LKC em no com
%  I1 e I2 entrando, bateria 12 V / 0,5 ohm, bateria com r=0,4 ohm,
%  no com cinco correntes, DOIS circuitos de duas malhas resolvidos
%  por Kirchhoff, quatro fontes ideais com terminais abertos e --- como
%  N4 --- DUAS BATERIAS DE FEM DIFERENTES EM PARALELO. Este ultimo tema
%  motivou a remocao de uma N4 do banco antigo, que o duplicava.
%  O banco de circuitos serie revisado ja traz regulacao de tensao e a
%  determinacao de fem e r por dois pontos de carga.
%  Este banco explora: fonte ABSORVENDO potencia, banco de baterias
%  serie x paralelo, LKC generalizada, contagem de equacoes
%  independentes, propagacao de incerteza na verificacao da LKC e o
%  teorema de Tellegen em forma elementar.
%  Distribuicao mantida: 6 N1 + 9 N2 + 9 N3 + 7 N4 = 31
% =====================================================================

\subsubsection*{Banco complementar --- Fontes reais e leis de Kirchhoff}

\begin{atencao}[Organização do banco]
A fonte real é o modelo $\EMF$ em série com $r$, e dele decorre
$V=\EMF-rI$ quando ela \emph{fornece} --- mas $V=\EMF+rI$ quando ela
\textbf{recebe} corrente, situação de carregamento que aparece com frequência e
costuma ser esquecida. As leis de Kirchhoff, por sua vez, decorrem da
conservação de carga e de energia: elas valem para \emph{qualquer} elemento,
linear ou não. O N3 explora fontes absorvendo potência, bancos de baterias e
contagem de equações; o N4 exige demonstração, ensaio e propagação de
incertezas.
\end{atencao}

\blocofixacao

\begin{questao}[1]{}
Uma fonte de $\EMF=\SI{6}{\volt}$ e $r=\SI{1}{\ohm}$ fornece
$\SI{1}{\ampere}$. Determine a tensão em seus terminais.
\resposta{$V=\SI{5}{\volt}$}
\resolucao{
\[
V=\EMF-rI=6-1(1)=\SI{5}{\volt}.
\]
A diferença de $\SI{1}{\volt}$ é a queda \emph{interna}, que não chega à carga.
Ela cresce com a corrente --- é por isso que a tensão de uma pilha cai quando se
liga um aparelho que exige mais.}
\end{questao}

\begin{questao}[1]{}
Uma fonte de $\EMF=\SI{9}{\volt}$ apresenta $\SI{8.4}{\volt}$ nos terminais
quando fornece $\SI{2}{\ampere}$. Determine sua resistência interna.
\resposta{$r=\SI{0.3}{\ohm}$}
\resolucao{
\[
r=\frac{\EMF-V}{I}=\frac{9-8{,}4}{2}=\SI{0.3}{\ohm}.
\]
\textbf{Método geral:} a queda interna dividida pela corrente. Uma única medição
sob carga, somada ao valor de $\EMF$, determina $r$ --- e $\EMF$ pode ser lida
em vazio com um voltímetro de alta impedância.}
\end{questao}

\begin{questao}[1]{}
Em um nó entram $\SI{2.5}{\ampere}$ e $\SI{1.5}{\ampere}$, e sai uma única
corrente. Determine seu valor.
\resposta{$I=\SI{4}{\ampere}$}
\resolucao{Pela lei dos nós,
\[
\sum I_{entra}=\sum I_{sai}
\;\Longrightarrow\;
2{,}5+1{,}5=\SI{4}{\ampere}.
\]
\textbf{Fundamento:} a LKC expressa a \textbf{conservação da carga}. Um nó não
acumula carga --- tudo o que entra, sai. Isso vale independentemente dos
elementos ligados a ele, e mesmo que sejam não lineares.}
\end{questao}

\begin{questao}[1]{}
Uma malha simples contém uma fonte ideal de $\SI{12}{\volt}$ e resistores de
$\SI{4}{\ohm}$ e $\SI{8}{\ohm}$. Determine a corrente pela lei das malhas.
\resposta{$I=\SI{1}{\ampere}$}
\resolucao{Pela LKT, a soma das quedas iguala a f.e.m.:
\[
12=4I+8I=12I
\;\Longrightarrow\;
I=\SI{1}{\ampere}.
\]
\textbf{Fundamento:} a LKT expressa a \textbf{conservação da energia}. Uma carga
que percorre a malha e retorna ao ponto de partida deve ter ganho e perdido a
mesma energia --- o potencial é uma função de estado.}
\end{questao}

\begin{questao}[1]{}
Sobre a diferença entre fonte ideal e fonte real, é correto afirmar:
\begin{alternativas}[2]
\item a ideal mantém a tensão constante para qualquer corrente
\item a real mantém a tensão constante para qualquer corrente
\item ambas mantêm a corrente constante
\item a ideal tem resistência interna maior
\end{alternativas}
\resposta{(a)}
\resolucao{A fonte \textbf{ideal} tem $r=0$: sua tensão terminal é
$\EMF$ independentemente da corrente, inclusive em curto-circuito --- o que
exigiria corrente infinita, revelando que se trata de uma idealização.\\
A fonte \textbf{real} tem $r>0$, e sua tensão cai como $\EMF-rI$.\\
\textbf{Quando a idealização é aceitável:} sempre que $r\ll R_{carga}$. Uma
fonte de bancada com $r=\SI{10}{\milli\ohm}$ alimentando
$\SI{100}{\ohm}$ é, para todos os efeitos, ideal.}
\end{questao}

\begin{questao}[1]{}
Uma pilha de $\EMF=\SI{1.5}{\volt}$ tem $r=\SI{0.25}{\ohm}$. Determine sua
corrente de curto-circuito.
\resposta{$I_{cc}=\SI{6}{\ampere}$}
\resolucao{
\[
I_{cc}=\frac{\EMF}{r}=\frac{1{,}5}{0{,}25}=\SI{6}{\ampere}.
\]
\textbf{Ressalva importante:} o valor pressupõe $r$ constante --- hipótese
otimista. Em uma pilha real, $r$ cresce com a corrente e com o aquecimento, e a
corrente efetiva é menor. Ainda assim, $\SI{6}{\ampere}$ em uma pilha pequena
significa $\SI{9}{\watt}$ dissipados \emph{dentro} dela: aquecimento rápido,
vazamento de eletrólito e risco de ruptura.}
\end{questao}

\blocoaplicacao

\begin{questao}[2]{}
Uma bateria de $\EMF=\SI{12}{\volt}$ e $r=\SI{0.6}{\ohm}$ alimenta
$R=\SI{5.4}{\ohm}$.
\begin{itensq}
\item Determine a corrente e a tensão terminal.
\item Determine as potências e o rendimento.
\end{itensq}
\resposta{$I=\SI{2}{\ampere}$, $V=\SI{10.8}{\volt}$;
$P_R=\SI{21.6}{\watt}$, $P_r=\SI{2.4}{\watt}$, $\eta=90\%$}
\resolucao{(a)
\[
I=\frac{\EMF}{r+R}=\frac{12}{6}=\SI{2}{\ampere};
\qquad
V=12-0{,}6(2)=\SI{10.8}{\volt}.
\]
(b)
\[
P_{total}=\EMF\,I=\SI{24}{\watt};
\quad
P_R=5{,}4(4)=\SI{21.6}{\watt};
\quad
P_r=0{,}6(4)=\SI{2.4}{\watt};
\]
\[
\eta=\frac{R}{R+r}=\frac{5{,}4}{6}=90\%.
\]
\textbf{Fórmula útil:} o rendimento de uma fonte real é sempre
$R/(R+r)$ --- ele depende só da razão entre as resistências, e não da tensão
nem da corrente.}
\end{questao}

\begin{questao}[2]{}
Um carregador de $\SI{14}{\volt}$ é ligado a uma bateria de
$\SI{12}{\volt}$, com resistência total de $\SI{0.4}{\ohm}$ no circuito.
\begin{itensq}
\item Determine a corrente de carga.
\item Determine a tensão nos terminais da bateria.
\end{itensq}
\resposta{$I=\SI{5}{\ampere}$; $V=\SI{14}{\volt}$}
\resolucao{(a) As duas f.e.m. estão em \textbf{oposição}, e a diferença é que
impulsiona a corrente:
\[
I=\frac{14-12}{0{,}4}=\SI{5}{\ampere}.
\]
(b) A bateria está \textbf{recebendo} corrente, então a queda interna
\emph{soma-se} à f.e.m.:
\[
V=\EMF+rI=12+0{,}4(5)=\SI{14}{\volt}.
\]
\textbf{Inversão do sinal:} quando a fonte fornece, $V=\EMF-rI<\EMF$; quando
recebe, $V=\EMF+rI>\EMF$. É por isso que a tensão de uma bateria de automóvel
sobe de $\SI{12.6}{\volt}$ para cerca de $\SI{14}{\volt}$ com o motor ligado ---
sinal de que o alternador está carregando.}
\end{questao}

\begin{questao}[2]{}
Em um nó, entram $\SI{3}{\ampere}$, $\SI{4}{\ampere}$ e $\SI{2}{\ampere}$;
sai $\SI{6}{\ampere}$ por um ramo e $I_5$ por outro. Determine $I_5$.
\resposta{$I_5=\SI{3}{\ampere}$, saindo}
\resolucao{
\[
3+4+2=6+I_5
\;\Longrightarrow\;
I_5=9-6=\SI{3}{\ampere}.
\]
O sinal positivo confirma o sentido admitido (saindo).\\
\textbf{Convenção alternativa:} adotar todas as correntes como \emph{entrando} e
impor soma nula. As que realmente saem aparecem com sinal negativo. As duas
convenções são equivalentes --- escolha uma e mantenha-a.}
\end{questao}

\begin{questao}[2]{}
Uma malha contém fontes ideais de $\SI{24}{\volt}$ e $\SI{6}{\volt}$ em
oposição, com resistores somando $\SI{6}{\ohm}$. Determine a corrente e
identifique qual fonte absorve energia.
\resposta{$I=\SI{3}{\ampere}$; a de $\SI{6}{\volt}$ absorve}
\resolucao{
\[
I=\frac{24-6}{6}=\SI{3}{\ampere},
\]
no sentido imposto pela fonte maior.\\
\textbf{Balanço:} a de $\SI{24}{\volt}$ fornece $24(3)=\SI{72}{\watt}$; a de
$\SI{6}{\volt}$ \emph{absorve} $6(3)=\SI{18}{\watt}$; os resistores dissipam
$6(9)=\SI{54}{\watt}$. Soma: $18+54=\SI{72}{\watt}$ \checkmark.\\
\textbf{Critério para identificar:} a fonte \emph{absorve} quando a corrente
entra pelo seu terminal positivo. Se ela for recarregável, está sendo
carregada; se não for (uma pilha comum), a energia vira calor e o dispositivo
pode ser danificado.}
\end{questao}

\begin{questao}[2]{}
Duas pilhas de $\SI{1.5}{\volt}$ e $r=\SI{0.2}{\ohm}$ cada são associadas.
Determine a f.e.m. e a resistência interna equivalentes em série e em paralelo.
\resposta{Série: $\SI{3}{\volt}$ e $\SI{0.4}{\ohm}$; paralelo:
$\SI{1.5}{\volt}$ e $\SI{0.1}{\ohm}$}
\resolucao{\textbf{Série:} as f.e.m. se somam e as resistências internas também:
$\EMF=\SI{3}{\volt}$, $r=\SI{0.4}{\ohm}$.\\
\textbf{Paralelo} (pilhas \emph{idênticas}): a f.e.m. permanece
$\SI{1.5}{\volt}$ e as resistências internas ficam em paralelo:
$r=\SI{0.1}{\ohm}$.\\
\textbf{Correntes de curto:} série $3/0{,}4=\SI{7.5}{\ampere}$; paralelo
$1{,}5/0{,}1=\SI{15}{\ampere}$.\\
\textbf{Advertência:} a associação em paralelo só é segura com pilhas
\emph{idênticas e no mesmo estado de carga}. Com f.e.m. diferentes, circula
corrente entre elas mesmo sem carga externa.}
\end{questao}

\begin{questao}[2]{}
Determine o número de equações independentes de LKC e de LKT para um circuito
com $n=5$ nós e $b=8$ ramos.
\resposta{$4$ de LKC e $4$ de LKT, totalizando $8$}
\resolucao{\textbf{LKC:} $n-1=4$ equações. A do quinto nó é combinação linear
das anteriores --- redundante.\\
\textbf{LKT:} $b-n+1=8-5+1=4$ equações, uma por malha independente.\\
\textbf{Total:} $4+4=8=b$ --- exatamente o número de correntes de ramo a
determinar. O sistema é quadrado e bem posto.\\
\textbf{Não é coincidência:} $(n-1)+(b-n+1)=b$ sempre. As leis de Kirchhoff
fornecem \emph{precisamente} as equações necessárias, nem uma a mais nem a
menos.}
\end{questao}

\begin{questao}[2]{}
Uma fonte de $\EMF=\SI{12}{\volt}$ e $r=\SI{0.5}{\ohm}$ alimenta duas cargas
em paralelo, de $\SI{12}{\ohm}$ e $\SI{6}{\ohm}$. Determine a tensão terminal e
a corrente de cada carga.
\resposta{$V=\SI{10.67}{\volt}$; $\SI{0.89}{\ampere}$ e
$\SI{1.78}{\ampere}$}
\resolucao{
\[
R_{par}=\frac{12(6)}{18}=\SI{4}{\ohm};
\qquad
I=\frac{12}{4{,}5}=\SI{2.67}{\ampere};
\qquad
V=12-0{,}5(2{,}67)=\SI{10.67}{\volt}.
\]
\emph{(Conferindo: $V=R_{par}I=4(2{,}67)=\SI{10.67}{\volt}$ \checkmark.)}
\[
I_{12}=\frac{10{,}67}{12}=\SI{0.89}{\ampere};
\qquad
I_{6}=\frac{10{,}67}{6}=\SI{1.78}{\ampere}.
\]
\textbf{Acoplamento pelas cargas:} se a carga de $\SI{6}{\ohm}$ for desligada,
a corrente total cai para $12/12{,}5=\SI{0.96}{\ampere}$ e a tensão sobe para
$\SI{11.52}{\volt}$ --- o que \emph{aumenta} a corrente da carga remanescente.
As cargas se influenciam através de $r$.}
\end{questao}

\begin{questao}[2]{}
Uma bateria em uso apresenta $\EMF=\SI{1.5}{\volt}$ e $r=\SI{0.3}{\ohm}$
quando nova, e $\EMF=\SI{1.3}{\volt}$ com $r=\SI{1.2}{\ohm}$ ao final da vida.
Compare a tensão entregue a uma carga de $\SI{3}{\ohm}$.
\resposta{$\SI{1.36}{\volt}$ contra $\SI{0.93}{\volt}$ --- queda de $32\%$}
\resolucao{\textbf{Nova:} $I=1{,}5/3{,}3=\SI{0.455}{\ampere}$ e
$V=3(0{,}455)=\SI{1.36}{\volt}$.\\
\textbf{Gasta:} $I=1{,}3/4{,}2=\SI{0.310}{\ampere}$ e
$V=3(0{,}310)=\SI{0.93}{\volt}$.\\
\textbf{Observe o que domina:} a f.e.m. caiu apenas $13\%$, mas a resistência
interna \textbf{quadruplicou} --- e é ela a responsável pela maior parte da
perda de desempenho.\\
\textbf{Consequência para diagnóstico:} medir a tensão de uma pilha \emph{em
vazio} é enganoso, pois ela mostra quase $\EMF$ mesmo já esgotada. O teste
correto é sob carga, ou a medição direta de $r$.}
\end{questao}

\begin{questao}[2]{}
Verifique o balanço de potência de um circuito com fonte de
$\EMF=\SI{20}{\volt}$, $r=\SI{1}{\ohm}$ e carga de $\SI{9}{\ohm}$.
\resposta{Fornecido $\SI{40}{\watt}$; dissipado $4+\SI{36}{\watt}$}
\resolucao{
\[
I=\frac{20}{10}=\SI{2}{\ampere};
\qquad
P_{fonte}=\EMF\,I=\SI{40}{\watt};
\]
\[
P_r=1(4)=\SI{4}{\watt};
\qquad
P_R=9(4)=\SI{36}{\watt};
\qquad 4+36=40\ \checkmark
\]
\textbf{Cuidado com a distinção:} a potência \emph{gerada} usa $\EMF$, não a
tensão terminal. Usando $V=\SI{18}{\volt}$ obter-se-ia $\SI{36}{\watt}$ ---
que é a potência \emph{entregue aos terminais}, e não a produzida internamente.
A diferença de $\SI{4}{\watt}$ é exatamente o que $r$ consome.}
\end{questao}

\blocoaprofundamento

\begin{questao}[3]{}
Duas fontes reais alimentam a mesma carga: $\EMF_1=\SI{12}{\volt}$ com
$r_1=\SI{1}{\ohm}$ e $\EMF_2=\SI{6}{\volt}$ com $r_2=\SI{1}{\ohm}$, ambas
ligadas a $R=\SI{4}{\ohm}$.
\begin{itensq}
\item Determine a tensão comum e as três correntes.
\item Identifique o comportamento de cada fonte.
\end{itensq}
\resposta{$V=\SI{8}{\volt}$; $I_1=\SI{4}{\ampere}$,
$I_2=\SI{-2}{\ampere}$, $I_R=\SI{2}{\ampere}$ --- a fonte 2 é
\textbf{carregada}}
\resolucao{(a) LKC no nó comum:
\[
\frac{V-12}{1}+\frac{V-6}{1}+\frac{V}{4}=0
\;\Longrightarrow\;
4V-48+4V-24+V=0
\;\Longrightarrow\;
V=\SI{8}{\volt}.
\]
\[
I_1=\frac{12-8}{1}=\SI{4}{\ampere};
\qquad
I_2=\frac{6-8}{1}=\SI{-2}{\ampere};
\qquad
I_R=\frac{8}{4}=\SI{2}{\ampere}.
\]
\textbf{LKC:} $4+(-2)=2$ \checkmark.\\
(b) \textbf{A fonte 2 tem corrente negativa} --- ela não fornece, \emph{recebe}
$\SI{2}{\ampere}$. A fonte 1, mais forte, alimenta a carga \emph{e} carrega a
fonte 2.\\
\textbf{Balanço completo:}
\begin{center}
\begin{tabular}{lr}
\hline
Fonte 1 fornece & $\SI{48}{\watt}$\\
$r_1$ dissipa & $\SI{16}{\watt}$\\
$R$ dissipa & $\SI{16}{\watt}$\\
Fonte 2 absorve (f.e.m.) & $\SI{12}{\watt}$\\
$r_2$ dissipa & $\SI{4}{\watt}$\\
\hline
\end{tabular}
\end{center}
Total consumido: $16+16+12+4=\SI{48}{\watt}$ \checkmark.\\
\textbf{Se a fonte 2 for recarregável}, os $\SI{12}{\watt}$ são armazenados. Se
não for, viram calor --- e o dispositivo aquece perigosamente.}
\end{questao}

\begin{questao}[3]{}
Um banco é formado por quatro pilhas de $\SI{1.5}{\volt}$ e
$r=\SI{0.2}{\ohm}$.
\begin{itensq}
\item Compare série e paralelo quanto a $\EMF$, $r$ e corrente de
      curto.
\item Determine a potência máxima disponível em cada arranjo.
\item Estabeleça o critério de escolha.
\end{itensq}
\resposta{Série: $\SI{6}{\volt}$/$\SI{0.8}{\ohm}$; paralelo:
$\SI{1.5}{\volt}$/$\SI{0.05}{\ohm}$; $P_{\max}=\SI{11.25}{\watt}$ nos
\textbf{dois}}
\resolucao{(a)
\begin{center}
\begin{tabular}{lccc}
\hline
Arranjo & $\EMF$ & $r$ & $I_{cc}$\\
\hline
Série & $\SI{6}{\volt}$ & $\SI{0.8}{\ohm}$ & $\SI{7.5}{\ampere}$\\
Paralelo & $\SI{1.5}{\volt}$ & $\SI{0.05}{\ohm}$ & $\SI{30}{\ampere}$\\
\hline
\end{tabular}
\end{center}
(b)
\[
\text{Série: } \frac{6^{2}}{4(0{,}8)}=\SI{11.25}{\watt};
\qquad
\text{Paralelo: } \frac{1{,}5^{2}}{4(0{,}05)}=\SI{11.25}{\watt}.
\]
\textbf{Idênticas} --- e não por acaso. Ambos os arranjos usam as mesmas quatro
pilhas, e a energia total disponível é a mesma. A topologia decide \emph{como}
ela é entregue, não \emph{quanto}.\\
(c) \textbf{Critério de escolha.}
\begin{itemize}
\item \textbf{Série} para cargas de \emph{alta} resistência, que precisam de
tensão. A carga ótima é $\SI{0.8}{\ohm}$.
\item \textbf{Paralelo} para cargas de \emph{baixa} resistência, que precisam de
corrente. A carga ótima é $\SI{0.05}{\ohm}$.
\item \textbf{Série-paralelo} (dois ramos de duas) daria $\SI{3}{\volt}$ e
$\SI{0.2}{\ohm}$ --- e, de novo, $\SI{11.25}{\watt}$.
\end{itemize}
\textbf{Regra:} escolha o arranjo cuja resistência interna se aproxime da carga
--- ou, mais realisticamente, cuja \emph{tensão} atenda ao que a carga exige.}
\end{questao}

\begin{questao}[3]{}
Aplique a LKC a uma \textbf{superfície fechada} que envolva dois nós de um
circuito, e explique por que o resultado é válido.
\begin{itensq}
\item Enuncie a forma generalizada.
\item Justifique a partir da conservação de carga.
\item Dê um exemplo de uso.
\end{itensq}
\resposta{A soma das correntes que atravessam qualquer superfície fechada é
nula}
\resolucao{(a) \textbf{LKC generalizada:} para qualquer superfície fechada que
corte o circuito, a soma algébrica das correntes que a atravessam é zero.\\
(b) \textbf{Justificativa:} a carga total dentro da superfície não pode variar
(em regime permanente). Como corrente é fluxo de carga, o que entra deve
igualar o que sai --- exatamente o mesmo argumento que vale para um nó, agora
aplicado a uma região que \emph{contém} vários nós.\\
Formalmente, escrevendo a LKC para cada nó interno e somando, as correntes dos
ramos \emph{internos} aparecem duas vezes com sinais opostos e se cancelam;
sobram apenas as que cruzam a fronteira.\\
(c) \textbf{Exemplo de uso.} Considere um bloco de circuito ligado ao restante
por três fios apenas. Se dois deles conduzem $\SI{5}{\ampere}$ e
$\SI{3}{\ampere}$ para dentro, o terceiro conduz necessariamente
$\SI{8}{\ampere}$ para fora --- \textbf{sem que seja preciso conhecer nada do
interior do bloco}, que pode conter dezenas de componentes.\\
\textbf{Aplicação prática:} é o que permite verificar a corrente de um circuito
inteiro medindo apenas os condutores de entrada, e é também o princípio do
\emph{dispositivo diferencial residual}, que compara fase e neutro --- qualquer
diferença indica corrente escapando por um terceiro caminho.}
\end{questao}

\begin{questao}[3]{}
Uma fonte foi ensaiada sob diferentes cargas, obtendo-se:
\begin{center}
\begin{tabular}{cccc}
\hline
$I$ ($\si{\ampere}$) & $0{,}50$ & $1{,}00$ & $2{,}00$\\
$V$ ($\si{\volt}$) & $11{,}75$ & $11{,}50$ & $11{,}00$\\
\hline
\end{tabular}
\end{center}
\begin{itensq}
\item Determine $\EMF$ e $r$.
\item Verifique a linearidade do modelo.
\item Estime a corrente de curto e discuta a extrapolação.
\end{itensq}
\resposta{$\EMF=\SI{12}{\volt}$, $r=\SI{0.5}{\ohm}$;
$I_{cc}=\SI{24}{\ampere}$}
\resolucao{(a) O modelo $V=\EMF-rI$ é uma reta. De dois pontos quaisquer:
\[
r=-\frac{\Delta V}{\Delta I}=-\frac{11{,}00-11{,}75}{2{,}00-0{,}50}
=\frac{0{,}75}{1{,}50}=\SI{0.5}{\ohm};
\]
\[
\EMF=11{,}75+0{,}5(0{,}50)=\SI{12}{\volt}.
\]
(b) \textbf{Verificação com o terceiro ponto:}
$12-0{,}5(1{,}00)=\SI{11.50}{\volt}$ \checkmark. Os três pontos são
\emph{exatamente} colineares --- o modelo linear descreve bem a fonte nessa
faixa.\\
\textbf{Nota de método:} com apenas dois pontos, sempre passa uma reta e a
linearidade não pode ser testada. O terceiro ponto é o que transforma o ajuste
em \emph{verificação}.\\
(c) $I_{cc}=\EMF/r=12/0{,}5=\SI{24}{\ampere}$.\\
\textbf{Cautela com a extrapolação:} o ensaio cobriu de $0{,}5$ a
$\SI{2}{\ampere}$, e o valor de curto está \emph{doze vezes} além do maior
ponto medido. Fontes reais aquecem, $r$ cresce e proteções atuam --- os
$\SI{24}{\ampere}$ são um \textbf{limite superior}, útil para dimensionar
proteção, não uma previsão.}
\end{questao}

\begin{questao}[3]{}
Um circuito tem $n=6$ nós, $b=9$ ramos e $2$ fontes de corrente ideais.
\begin{itensq}
\item Determine o número de equações de Kirchhoff necessárias.
\item Explique como as fontes de corrente reduzem o trabalho.
\item Compare com a contagem para análise nodal.
\end{itensq}
\resposta{$5$ LKC $+$ $4$ LKT $=9$; as fontes de corrente fixam duas correntes
de ramo, restando $7$ incógnitas}
\resolucao{(a) LKC: $n-1=5$. LKT: $b-n+1=9-6+1=4$. Total: $9$ --- igual ao
número de correntes de ramo.\\
(b) \textbf{Cada fonte de corrente ideal fixa diretamente a corrente de seu
ramo.} Duas delas eliminam duas incógnitas, restando $7$. Em compensação, a
tensão sobre cada fonte de corrente é desconhecida, o que impede escrever a LKT
de malhas que a atravessem --- é preciso escolher malhas que as contornem, ou
tratar essas tensões como novas incógnitas.\\
\textbf{Resultado líquido:} sistema $7\times7$ em vez de $9\times9$.\\
(c) \textbf{Análise nodal:} $n-1=5$ equações, com as fontes de corrente entrando
\emph{diretamente} no vetor independente, sem nenhuma manipulação. Sistema
$5\times5$.\\
\textbf{Conclusão:} a análise nodal é claramente superior aqui. Ela é, na
verdade, uma forma organizada de aplicar a LKC com a LKT já embutida na escolha
dos potenciais --- e por isso produz o sistema menor sempre que há fontes de
corrente.}
\end{questao}

\begin{questao}[3]{}
Uma fonte tem resistência interna que \emph{cresce} com a corrente, segundo
$r=0{,}2+0{,}1\,I$ (com $r$ em $\si{\ohm}$ e $I$ em $\si{\ampere}$), com
$\EMF=\SI{6}{\volt}$.
\begin{itensq}
\item Determine a corrente para uma carga de $\SI{1}{\ohm}$.
\item Compare com a previsão do modelo linear com $r=\SI{0.2}{\ohm}$.
\end{itensq}
\resposta{$I=\SI{3.80}{\ampere}$ contra $\SI{5}{\ampere}$ previstos}
\resolucao{(a) A LKT dá
\[
6=(0{,}2+0{,}1I)I+1\cdot I
\;\Longrightarrow\;
0{,}1I^{2}+1{,}2I-6=0.
\]
\[
I=\frac{-1{,}2+\sqrt{1{,}44+2{,}4}}{0{,}2}
=\frac{-1{,}2+1{,}960}{0{,}2}=\SI{3.80}{\ampere}.
\]
\emph{(Verificação: $r=0{,}2+0{,}380=\SI{0.58}{\ohm}$ e
$6/(0{,}58+1)=\SI{3.80}{\ampere}$ \checkmark.)}\\
(b) \textbf{Modelo linear:} $I=6/1{,}2=\SI{5}{\ampere}$ --- \textbf{$32\%$
acima} do valor real.\\
\textbf{Duas lições.}\\
\emph{(i)} As \textbf{leis de Kirchhoff continuam válidas} --- elas não exigem
linearidade. O que muda é que a equação resultante deixa de ser linear e passa a
exigir solução de segundo grau.\\
\emph{(ii)} O modelo de $r$ constante é uma \emph{aproximação de pequenos
sinais}. Ele descreve bem a fonte perto do ponto onde $r$ foi medido, e falha
progressivamente longe dele. Baterias reais comportam-se assim, e é por isso
que fabricantes especificam $r$ para uma corrente de ensaio definida.}
\end{questao}

\begin{questao}[3]{}
Uma bateria de $\EMF=\SI{12}{\volt}$ e $r=\SI{0.5}{\ohm}$ deve entregar tensão
terminal de pelo menos $\SI{11.4}{\volt}$.
\begin{itensq}
\item Determine a corrente máxima admissível.
\item Determine a menor carga que pode ser conectada.
\item Determine a potência entregue nessa condição.
\end{itensq}
\resposta{$I\le\SI{1.2}{\ampere}$; $R\ge\SI{9.5}{\ohm}$;
$P=\SI{13.68}{\watt}$}
\resolucao{(a) De $V=\EMF-rI\ge11{,}4$:
\[
0{,}5I\le0{,}6
\;\Longrightarrow\;
I\le\SI{1.2}{\ampere}.
\]
(b) $R=V/I=11{,}4/1{,}2=\SI{9.5}{\ohm}$ --- cargas menores exigiriam mais
corrente e derrubariam a tensão.\\
(c) $P=VI=11{,}4(1{,}2)=\SI{13.68}{\watt}$, com rendimento
$9{,}5/10=95\%$.\\
\textbf{Contraste com a máxima transferência:} a potência máxima ocorreria em
$R=\SI{0.5}{\ohm}$, valendo $\SI{72}{\watt}$ --- mas com tensão terminal de
apenas $\SI{6}{\volt}$, muito abaixo do requisito.\\
\textbf{A restrição de tensão é quase sempre a que manda} em eletrônica, porque
circuitos alimentados simplesmente não funcionam abaixo de sua tensão mínima. A
saída de engenharia é reduzir $r$, não aceitar tensão menor.}
\end{questao}

\begin{questao}[3]{}
Um circuito contém uma fonte \textbf{dependente} de tensão de valor
$3\,I_x$ (em volts, com $I_x$ em ampères) em série com $\SI{2}{\ohm}$, além de
uma fonte independente de $\SI{20}{\volt}$ e um resistor de $\SI{8}{\ohm}$, em
malha única. A corrente de controle $I_x$ é a própria corrente da malha.
\begin{itensq}
\item Escreva a LKT e resolva.
\item Explique por que Kirchhoff continua válido.
\end{itensq}
\resposta{$I=\SI{1.54}{\ampere}$ (fonte dependente em oposição)}
\resolucao{(a) Percorrendo a malha, com a fonte dependente \emph{opondo-se}:
\[
20=2I+8I+3I=13I
\;\Longrightarrow\;
I=\frac{20}{13}=\SI{1.54}{\ampere}.
\]
\emph{(Se a fonte dependente auxiliasse, viria $20=10I-3I=7I$ e
$I=\SI{2.86}{\ampere}$ --- o sentido da fonte dependente altera
completamente o resultado, e por isso a convenção de polaridade precisa estar
explícita no esquema.)}\\
(b) \textbf{As leis de Kirchhoff decorrem de conservação de carga e de
energia}, e nenhuma dessas conservações depende de os elementos serem lineares
ou independentes. A LKT continua sendo "a soma das quedas em uma malha é nula";
a LKC continua sendo "o que entra em um nó, sai".\\
\textbf{O que muda:} a fonte dependente introduz uma \textbf{equação de
restrição} adicional, ligando seu valor a uma variável do circuito. Essa
equação é resolvida \emph{junto} com as de Kirchhoff.\\
\textbf{Consequência estrutural:} o termo migra para o lado das incógnitas, o
que pode tornar a matriz assimétrica --- como se observa na análise nodal e na
de malhas. A reciprocidade se perde; as leis, não.}
\end{questao}

\begin{questao}[3]{}
Uma fonte alimenta uma carga através de um cabo de $\SI{0.3}{\ohm}$ (ida e
volta). A fonte tem $\EMF=\SI{24}{\volt}$ e $r=\SI{0.2}{\ohm}$; a carga vale
$\SI{9.5}{\ohm}$.
\begin{itensq}
\item Determine a corrente, a tensão na carga e a tensão nos terminais da
      fonte.
\item Determine onde a energia é perdida.
\end{itensq}
\resposta{$I=\SI{2.4}{\ampere}$; $V_{carga}=\SI{22.8}{\volt}$;
$V_{fonte}=\SI{23.52}{\volt}$}
\resolucao{(a)
\[
I=\frac{24}{0{,}2+0{,}3+9{,}5}=\frac{24}{10}=\SI{2.4}{\ampere};
\]
\[
V_{carga}=9{,}5(2{,}4)=\SI{22.8}{\volt};
\qquad
V_{fonte}=24-0{,}2(2{,}4)=\SI{23.52}{\volt}.
\]
A diferença de $\SI{0.72}{\volt}$ entre os dois é exatamente a queda no cabo.\\
(b) \textbf{Distribuição das perdas:}
\begin{center}
\begin{tabular}{lrr}
\hline
Elemento & Potência & Fração\\
\hline
Carga ($\SI{9.5}{\ohm}$) & $\SI{54.72}{\watt}$ & $95{,}0\%$\\
Cabo ($\SI{0.3}{\ohm}$) & $\SI{1.73}{\watt}$ & $3{,}0\%$\\
Interna ($\SI{0.2}{\ohm}$) & $\SI{1.15}{\watt}$ & $2{,}0\%$\\
\hline
Total & $\SI{57.6}{\watt}$ & $100\%$\\
\hline
\end{tabular}
\end{center}
\textbf{Observação de diagnóstico:} medindo apenas nos terminais da fonte
($\SI{23.52}{\volt}$), a instalação pareceria melhor do que é. A tensão que
importa é a que chega à \emph{carga} --- e ela está $5\%$ abaixo da f.e.m.\\
\textbf{Regra de medição:} sempre meça no ponto de uso, não na fonte. As duas
leituras diferem exatamente pela queda no caminho.}
\end{questao}

\blocodesafio

\begin{questao}[4]{}
\textbf{(Ensaio de caracterização.)} Descreva como obter $\EMF$ e $r$ de uma
fonte por ensaios em vazio e em curto-circuito.
\begin{itensq}
\item Deduza as relações.
\item Analise os riscos do ensaio em curto.
\item Proponha uma alternativa segura e compare.
\end{itensq}
\resposta{$\EMF=V_{oc}$ e $r=V_{oc}/I_{cc}$; o curto é destrutivo em fontes de
baixa impedância}
\resolucao{(a) \textbf{Em vazio} ($I=0$): $V_{oc}=\EMF$ --- a leitura direta,
desde que o voltímetro tenha impedância muito maior que $r$.\\
\textbf{Em curto} ($V=0$): $\EMF=rI_{cc}$, logo
\[
r=\frac{V_{oc}}{I_{cc}}.
\]
(b) \textbf{Riscos do curto.}
\begin{itemize}
\item Uma pilha AA de $\SI{1.5}{\volt}$ com $r=\SI{0.25}{\ohm}$ entrega
$\SI{6}{\ampere}$ e dissipa $\SI{9}{\watt}$ internamente --- aquecimento
rápido e vazamento.
\item Uma bateria de automóvel ($r\approx\SI{5}{\milli\ohm}$) chega a
$\SI{2500}{\ampere}$, com centelhamento, fusão de terminais e risco de
explosão por ignição do hidrogênio liberado.
\item Baterias de lítio podem entrar em fuga térmica e incendiar.
\item Além disso, $r$ \emph{aumenta} durante o curto, de modo que a leitura
subestima a corrente inicial e superestima $r$.
\end{itemize}
(c) \textbf{Alternativa segura: dois pontos de carga.} Meça $(I_1,V_1)$ e
$(I_2,V_2)$ com resistências conhecidas e moderadas, e ajuste a reta:
\[
r=-\frac{V_2-V_1}{I_2-I_1};
\qquad
\EMF=V_1+rI_1.
\]
\textbf{Comparação.} O método de dois pontos é seguro, não perturba a fonte, e
com três ou mais pontos ainda \emph{testa} a linearidade do modelo --- coisa que
o ensaio de curto não faz. O ensaio em curto só se justifica em fontes
projetadas para suportá-lo, com limitação eletrônica de corrente.}
\end{questao}

\begin{questao}[4]{}
\textbf{(Demonstração --- balanço de potência.)} Prove que a soma algébrica das
potências de todos os elementos de um circuito é nula.
\begin{itensq}
\item Estabeleça a convenção de sinais.
\item Demonstre a partir das leis de Kirchhoff.
\item Explique por que o resultado independe da natureza dos elementos.
\end{itensq}
\resposta{$\sum_k v_ki_k=0$ --- consequência direta de LKC e LKT}
\resolucao{(a) \textbf{Convenção do receptor:} para cada ramo $k$, adote a
corrente $i_k$ entrando pelo terminal de maior potencial. A potência
$p_k=v_ki_k$ é então \emph{positiva} se o elemento absorve e \emph{negativa} se
fornece.\\
(b) \textbf{Demonstração.} Escreva cada tensão de ramo como diferença de
potenciais nodais: $v_k=V_a-V_b$, onde $a$ e $b$ são os nós do ramo $k$. Então
\[
\sum_k p_k=\sum_k (V_a-V_b)\,i_k.
\]
Reagrupando por \emph{nó} em vez de por ramo, cada potencial $V_n$ aparece
multiplicado pela soma algébrica das correntes que chegam ao nó $n$:
\[
\sum_k p_k=\sum_n V_n\left(\sum_{k\in n}\pm i_k\right)=\sum_n V_n\cdot 0=0,
\]
pois o parêntese é exatamente a \textbf{LKC} do nó $n$.\\
\emph{(A escrita $v_k=V_a-V_b$ já embute a \textbf{LKT}: ela pressupõe que o
potencial é bem definido em cada nó.)}\\
(c) \textbf{Independência da natureza dos elementos.} A demonstração usou
\emph{apenas} as duas leis de Kirchhoff e a definição $p=vi$. Em nenhum momento
foi preciso saber a relação entre $v_k$ e $i_k$ de cada elemento.\\
Portanto o resultado vale para resistores, fontes, diodos, capacitores,
indutores, elementos ativos --- lineares ou não, com ou sem memória. É o
\textbf{teorema de Tellegen}, e sua generalidade é notável: ele é uma
propriedade da \emph{topologia} do circuito, não dos componentes.\\
\textbf{Uso prático:} verificar o balanço de potência é o teste de consistência
mais completo de qualquer solução --- se ele não fecha, há erro em algum lugar.}
\end{questao}

\begin{questao}[4]{}
\textbf{(Formulação com fonte ideal entre dois nós.)} Um circuito tem dois nós
não referenciais, $A$ e $B$, com uma fonte ideal de tensão $\EMF$ ligada
diretamente entre eles.
\begin{itensq}
\item Explique por que a LKC não pode ser escrita isoladamente em $A$ nem em
      $B$.
\item Formule o sistema corretamente.
\item Determine a corrente na fonte após resolver.
\end{itensq}
\resposta{Usar supernó: uma LKC envolvendo $A$ e $B$, mais a restrição
$V_A-V_B=\EMF$}
\resolucao{(a) A corrente que atravessa uma fonte de tensão \textbf{ideal} é
determinada pelo circuito, não pela fonte --- não existe relação
$i=f(v)$ para ela. Ao escrever a LKC em $A$, essa corrente desconhecida aparece
como uma incógnita a mais, e o mesmo ocorre em $B$. Duas equações, três
incógnitas.\\
(b) \textbf{Formulação por supernó.} Englobe $A$ e $B$ em uma única superfície
fechada. Pela LKC generalizada, a corrente interna da fonte \emph{se cancela}
--- ela entra em um nó e sai do outro:
\[
\sum_{\text{ramos que cruzam a fronteira}} i=0.
\]
Essa é \emph{uma} equação. A segunda vem da própria fonte:
\[
V_A-V_B=\EMF.
\]
Duas equações, duas incógnitas ($V_A$ e $V_B$). \textbf{Sistema fechado.}\\
(c) \textbf{Corrente na fonte:} obtida \emph{depois}, aplicando a LKC a
\emph{um} dos dois nós isoladamente. Com $V_A$ e $V_B$ já conhecidos, todas as
demais correntes daquele nó são calculáveis, e a da fonte é o que falta para
fechar o balanço.\\
\textbf{Observação:} se um dos terminais da fonte for o próprio nó de
referência, o problema desaparece --- o outro potencial fica \emph{conhecido} e
deixa de ser incógnita. \textbf{Escolher o terra em um terminal de fonte é a
simplificação mais barata que existe.}}
\end{questao}

\begin{questao}[4]{}
\textbf{(Projeto de verificação experimental.)} Elabore um ensaio de laboratório
para verificar a LKC em um nó com três ramos, considerando incertezas dos
instrumentos.
\begin{itensq}
\item Descreva o procedimento.
\item Propague as incertezas e estabeleça o critério de aceitação.
\item Interprete um resíduo de $\SI{0.35}{\ampere}$ em correntes de alguns
      ampères.
\end{itensq}
\resposta{Critério: resíduo compatível com $u_c\approx\sqrt{\sum u_k^{2}}$;
$\SI{0.35}{\ampere}$ indica erro real}
\resolucao{(a) \textbf{Procedimento.} Meça as três correntes
\emph{simultaneamente}, com três amperímetros --- ou, se houver apenas um,
faça as três medições sem alterar o circuito entre elas. Registre sentidos e
adote uma convenção única (por exemplo, positivo entrando).\\
\textbf{Cuidado essencial:} inserir o amperímetro \emph{altera} o circuito, pois
acrescenta sua resistência interna ao ramo. Com três inserções sucessivas, o
circuito muda três vezes e a verificação perde sentido. Use amperímetro de baixa
resistência ou, melhor, alicate amperímetro, que não abre o ramo.\\
(b) \textbf{Propagação.} Suponha
$I_1=(5{,}02\pm0{,}05)\,\si{\ampere}$,
$I_2=(3{,}01\pm0{,}03)\,\si{\ampere}$ entrando, e $I_3$ saindo. A previsão é
$I_1+I_2=\SI{8.03}{\ampere}$, com incerteza combinada
\[
u_c=\sqrt{0{,}05^{2}+0{,}03^{2}}=\SI{0.058}{\ampere}.
\]
Incluindo a incerteza da própria medição de $I_3$ (digamos
$\SI{0.08}{\ampere}$), o resíduo esperado tem incerteza
$\sqrt{0{,}058^{2}+0{,}08^{2}}=\SI{0.099}{\ampere}$.\\
\textbf{Critério de aceitação:} $|I_1+I_2-I_3|\le2u\approx
\SI{0.20}{\ampere}$ (dois desvios, cerca de $95\%$ de confiança).\\
(c) \textbf{Resíduo de $\SI{0.35}{\ampere}$} é $3{,}5$ vezes a incerteza
combinada --- \textbf{não é compatível com erro aleatório}. Causas prováveis, em
ordem:
\begin{itemize}
\item \textbf{Erro de sinal ou de sentido} em uma das leituras --- o mais comum.
\item \textbf{Um quarto ramo não contabilizado}, inclusive uma fuga para o terra
ou pela bancada.
\item \textbf{Instrumento descalibrado} ou operando fora de escala.
\item \textbf{Circuito alterado entre medições} (aquecimento, contato
intermitente).
\end{itemize}
\textbf{Conclusão importante:} a LKC \emph{não} pode ser violada --- ela é
conservação de carga. Um resíduo significativo é sempre erro de medição ou de
modelo, nunca da lei. O ensaio, portanto, verifica o \emph{procedimento}, e não
a física.}
\end{questao}

\begin{questao}[4]{}
\textbf{(Baterias em série desequilibradas.)} Quatro células de
$\SI{1.5}{\volt}$ estão em série alimentando uma carga, mas uma delas está
gasta, com $\EMF=\SI{0.9}{\volt}$ e $r=\SI{1.5}{\ohm}$ (as demais mantêm
$r=\SI{0.2}{\ohm}$). A carga vale $\SI{5}{\ohm}$.
\begin{itensq}
\item Determine a corrente e a tensão em cada célula.
\item Identifique o que ocorre com a célula gasta.
\item Discuta a implicação prática.
\end{itensq}
\resposta{$I=\SI{0.76}{\ampere}$; a célula gasta apresenta
$\SI{-0.24}{\volt}$ --- está sendo \textbf{invertida}}
\resolucao{(a)
\[
\EMF_{total}=3(1{,}5)+0{,}9=\SI{5.4}{\volt};
\qquad
r_{total}=3(0{,}2)+1{,}5=\SI{2.1}{\ohm};
\]
\[
I=\frac{5{,}4}{2{,}1+5}=\frac{5{,}4}{7{,}1}=\SI{0.76}{\ampere}.
\]
\textbf{Tensão de cada célula boa:} $1{,}5-0{,}2(0{,}76)=\SI{1.35}{\volt}$.\\
\textbf{Tensão da célula gasta:}
$0{,}9-1{,}5(0{,}76)=0{,}9-1{,}14=\SI{-0.24}{\volt}$.\\
(b) \textbf{A tensão da célula gasta é negativa} --- ou seja, ela está sendo
percorrida por corrente no sentido de \emph{carga}, forçada pelas três células
boas. Ela deixou de ser fonte e passou a ser carga, dissipando
$1{,}5(0{,}76)^{2}=\SI{0.87}{\watt}$ internamente.\\
\textbf{Em células não recarregáveis}, isso provoca a chamada \textbf{inversão
de polaridade}: reações internas indesejadas, geração de gás, aquecimento e
vazamento de eletrólito --- exatamente o que se observa quando uma pilha
"vaza" em um controle remoto esquecido ligado.\\
(c) \textbf{Implicações práticas.}
\begin{itemize}
\item \textbf{Nunca misture pilhas novas e usadas}, nem de marcas ou tecnologias
diferentes, em série. A mais fraca será invertida pelas demais.
\item \textbf{Substitua sempre o conjunto completo.}
\item Em baterias de lítio em série, o problema é grave a ponto de exigir
\textbf{circuito de balanceamento} (BMS), que monitora cada célula
individualmente e impede que qualquer uma seja levada à inversão.
\end{itemize}
\textbf{Diagnóstico:} a célula gasta domina o conjunto por sua \emph{resistência
interna}, não por sua f.e.m. --- $\SI{1.5}{\ohm}$ contra
$\SI{0.6}{\ohm}$ das outras três somadas.}
\end{questao}

\begin{questao}[4]{}
\textbf{(Modelagem --- fonte com limitação de corrente.)} Uma fonte de bancada
tem $\EMF=\SI{20}{\volt}$ e $r=\SI{0.1}{\ohm}$, com limitação eletrônica em
$\SI{2}{\ampere}$.
\begin{itensq}
\item Descreva a característica $V\times I$ completa.
\item Determine o ponto de operação para cargas de $\SI{50}{\ohm}$,
      $\SI{10}{\ohm}$ e $\SI{2}{\ohm}$.
\item Explique por que o modelo de fonte real deixa de valer.
\end{itensq}
\resposta{Duas regiões: fonte de tensão até $\SI{2}{\ampere}$, fonte de corrente
depois}
\resolucao{(a) \textbf{Característica em dois trechos.}
\begin{itemize}
\item \textbf{Região de tensão constante} ($I<\SI{2}{\ampere}$):
$V=20-0{,}1I$ --- reta quase horizontal.
\item \textbf{Região de corrente constante} ($V<\SI{19.8}{\volt}$):
$I=\SI{2}{\ampere}$ --- reta vertical.
\end{itemize}
O "joelho" fica em $(\SI{2}{\ampere};\ \SI{19.8}{\volt})$, com
$R_{joelho}=\SI{9.9}{\ohm}$.\\
(b)
\begin{center}
\begin{tabular}{lccl}
\hline
$R_L$ & $I$ & $V$ & Região\\
\hline
$\SI{50}{\ohm}$ & $\SI{0.399}{\ampere}$ & $\SI{19.96}{\volt}$ & tensão\\
$\SI{10}{\ohm}$ & $\SI{1.98}{\ampere}$ & $\SI{19.80}{\volt}$ & tensão (limite)\\
$\SI{2}{\ohm}$ & $\SI{2.00}{\ampere}$ & $\SI{4.00}{\volt}$ & corrente\\
\hline
\end{tabular}
\end{center}
Note o terceiro caso: o modelo linear preveria
$20/2{,}1=\SI{9.5}{\ampere}$ --- quase cinco vezes o real.\\
(c) \textbf{Por que o modelo falha.} A fonte real é um modelo \emph{linear} de
dois parâmetros, e sua característica é uma \emph{única} reta. A limitação
eletrônica introduz um segundo trecho, tornando a característica
\textbf{quebrada} --- e nenhuma reta descreve as duas regiões.\\
\textbf{Como proceder corretamente:}
\begin{itemize}
\item Determinar em qual região a carga opera, comparando $R_L$ com
$R_{joelho}=\SI{9.9}{\ohm}$.
\item $R_L>R_{joelho}$: use o modelo de \emph{fonte de tensão} $(\EMF,r)$.
\item $R_L<R_{joelho}$: use o modelo de \emph{fonte de corrente} ideal de
$\SI{2}{\ampere}$, e a tensão sai de $V=R_LI$.
\end{itemize}
\textbf{As leis de Kirchhoff seguem válidas nas duas regiões} --- o que muda é
apenas a relação constitutiva da fonte. É o mesmo procedimento usado com
diodos: identificar o trecho e aplicar o modelo correspondente.}
\end{questao}

\begin{questao}[4]{}
\textbf{(Sistema completo por Kirchhoff.)} Um circuito de duas malhas tem
$\EMF_1=\SI{18}{\volt}$ com $r_1=\SI{1}{\ohm}$ no primeiro ramo,
$\EMF_2=\SI{12}{\volt}$ com $r_2=\SI{2}{\ohm}$ no segundo, e
$R=\SI{6}{\ohm}$ no ramo comum.
\begin{itensq}
\item Escreva as equações de Kirchhoff e resolva.
\item Verifique o balanço de potência.
\item Identifique o papel de cada fonte.
\end{itensq}
\resposta{$I_1=\SI{3.6}{\ampere}$, $I_2=\SI{-1.2}{\ampere}$,
$I_R=\SI{2.4}{\ampere}$}
\resolucao{(a) Chamando $V$ a tensão do nó comum:
\[
\frac{V-18}{1}+\frac{V-12}{2}+\frac{V}{6}=0.
\]
Multiplicando por $6$: $6V-108+3V-36+V=0\Rightarrow10V=144
\Rightarrow V=\SI{14.4}{\volt}$.
\[
I_1=\frac{18-14{,}4}{1}=\SI{3.6}{\ampere};
\quad
I_2=\frac{12-14{,}4}{2}=\SI{-1.2}{\ampere};
\quad
I_R=\frac{14{,}4}{6}=\SI{2.4}{\ampere}.
\]
\textbf{LKC:} $3{,}6-1{,}2=2{,}4$ \checkmark.\\
(b) \textbf{Balanço.}
\begin{center}
\begin{tabular}{lr}
\hline
$\EMF_1$ fornece $18(3{,}6)$ & $\SI{64.8}{\watt}$\\
$r_1$ dissipa $1(3{,}6)^{2}$ & $\SI{12.96}{\watt}$\\
$R$ dissipa $6(2{,}4)^{2}$ & $\SI{34.56}{\watt}$\\
$\EMF_2$ absorve $12(1{,}2)$ & $\SI{14.4}{\watt}$\\
$r_2$ dissipa $2(1{,}2)^{2}$ & $\SI{2.88}{\watt}$\\
\hline
Total consumido & $\SI{64.8}{\watt}$\\
\hline
\end{tabular}
\end{center}
Fecha exatamente \checkmark.\\
(c) \textbf{Papéis.} A fonte 1 é a única \emph{geradora}: ela alimenta a carga
$R$ \emph{e} carrega a fonte 2. A fonte 2 comporta-se como \textbf{receptor},
porque sua f.e.m. ($\SI{12}{\volt}$) é menor que a tensão do nó
($\SI{14.4}{\volt}$) --- e corrente sempre entra por onde o potencial externo é
maior.\\
\textbf{Critério geral:} compare a f.e.m. de cada fonte com a tensão do nó a
que está ligada. Maior $\Rightarrow$ fornece; menor $\Rightarrow$ recebe. O
sinal da corrente calculada confirma automaticamente.}
\end{questao}
